%% kapa: exam=1 round=37 year=2026 period=1 form=A subject=economics subject_ko=경제학원론 seq=3 number=43
%% source: Q-Net 감정평가사 기출문제 (cst00309) artlSeq=5258924 file=raw/1008/5258924/2026년 제37회 감정평가사 1차 1교시.pdf
%% conversion: kordoc 4.24.51 -> tools/gen1.py (auto); review=proofread-sonnet; flags=none
\begin{question}{43}
효용극대화를 추구하는 소비자 A가 X, Y 두 재화를 $MRS_{XY}<\dfrac{P_{X}}{P_{Y}}$인 점에서 소비하고 있을 때, 아래 설명 중 옳은 것을 모두 고른 것은? (단, 한계대체율 $(MRS_{XY})=\left|\dfrac{dY}{dX}\right|$이고 체감하며, 재화 X, Y의 가격은 각각 $P_{X}$, $P_{Y}$이고, X $>$ 0, Y $>$ 0이다.)
\begin{framedtext}
ㄱ. 효용극대화를 위해서는 X재를 더 소비하고, Y재를 덜 소비해야 한다. \par
ㄴ. $MRS_{XY}$는 소비자 A가 동일한 효용수준을 유지하면서 X재 1단위를 더 소비하기 위해 포기할 용의가 있는 Y재의 양이다. \par
ㄷ. $P_{X}=2$, $P_{Y}=1$일 때, 시장에서는 Y재 1단위와 X재 2단위가 교환된다. \par
ㄹ. X재에 1원을 지출했을 때 얻는 한계효용은 Y재에 1원을 지출했을 때 얻는 한계효용보다 작다. \par
\end{framedtext}
\begin{choices}
  \item ㄱ, ㄴ
  \item ㄱ, ㄷ
  \item ㄴ, ㄷ
  \item ㄴ, ㄹ
  \item ㄷ, ㄹ
\end{choices}
\end{question}
\answer{4}
